\documentclass[12pt, letterpaper]{article}

\title{CS162: Arrays and Loops}
\author{Amiel L. Iliesi}
\date{2026-03-27}

% packages
\usepackage[hidelinks]{hyperref}
\usepackage{ulem}

% commands
\newcommand{\question}[2][4em]{\textbf{#2}\vspace{#1}\par}

% global options
\setlength{\parindent}{0pt}

\begin{document}

\maketitle

\tableofcontents

\pagebreak
\section{Arrays}

\subsection{Basic Arrays}
\question{Initialize an array with the values \texttt{\{1, 2, 3, 4\}}.}

\question[8em]{Declare an integer array with room for 3 values. Then in individual
assignment statements, set each element to the first 3 square numbers.}

\subsection{Character Arrays}

\question{Set a character array to say \texttt{"Hello, World"}.}

\question[8em]{Declare a chracter array with enough room to hold the phrase
\texttt{"Hi"}. Then manually set each character to make a valid c-string.}

\pagebreak
\section{Loops}

\subsection{While Loops}
\question[6em]{Given a c-string with the name \texttt{str}, use a \uline{while} loop
to count the number of a's in the string.}

\subsection{Do-While Loops}
\question[8em]{Let the user enter sentences. Output them back to the user but in
quotes. Stop when they say they are finished. Use a \uline{do-while} loop.}

\subsection{For Loops}
\question{Given an array of integers named \texttt{ints}, and an
\texttt{unsigned int} denoting the array size named \texttt{ints\_size}, square
the elements, in-place, using a \uline{for} loop.}

\end{document}